% =====================================================================
%  Noxed / SprintLab — Research Questions, Student-Case Taxonomy,
%  Behavioural Metrics, and the Question-Difficulty Formula
%  Compile:  pdflatex Noxed_Research_Questions.tex   (run twice for refs)
% =====================================================================
\documentclass[11pt]{article}
\usepackage[a4paper,margin=2.2cm]{geometry}
\usepackage{lmodern}
\usepackage[T1]{fontenc}
\usepackage[utf8]{inputenc}
\usepackage{amsmath,amssymb}
\usepackage{booktabs}
\usepackage{array}
\usepackage{tabularx}
\usepackage{xcolor}
\usepackage{enumitem}
\usepackage{titlesec}
\usepackage{mdframed}
\usepackage[hidelinks]{hyperref}

\definecolor{nblue}{HTML}{2F4B7C}
\definecolor{ngreen}{HTML}{1B7F5B}
\definecolor{norange}{HTML}{B5651D}
\definecolor{npurple}{HTML}{7A3B8F}
\definecolor{ngray}{HTML}{4A4A4A}
\definecolor{lightbg}{HTML}{F4F6F9}

\titleformat{\section}{\normalfont\Large\bfseries\color{nblue}}{\thesection}{0.6em}{}
\titleformat{\subsection}{\normalfont\large\bfseries\color{ngray}}{\thesubsection}{0.6em}{}
\newcolumntype{Y}{>{\raggedright\arraybackslash}X}
\newcommand{\rqbox}[2]{\begin{mdframed}[backgroundcolor=lightbg,linecolor=nblue,linewidth=1pt,roundcorner=3pt,skipabove=4pt,skipbelow=4pt]\textbf{\color{nblue}#1}\, #2\end{mdframed}}
\setlist{itemsep=1pt,topsep=2pt}
\renewcommand{\arraystretch}{1.25}

\title{\vspace{-1.2cm}\textbf{\color{nblue}Noxed --- Adaptive Learning for Egyptian School Science}\\[2pt]
\large Research Questions, Student-Case Taxonomy, Behavioural Metrics\\ and a Question-Difficulty Formula}
\author{\normalsize SprintLab --- Research \& Design Note}
\date{\normalsize \today}

\begin{document}
\maketitle
\vspace{-0.6cm}

\begin{abstract}\noindent
\textbf{Noxed} is a competitive quiz game that teaches \textbf{science} to Egyptian school students and, after every match, recommends \emph{the next question / level} that will move the learner forward the most. It is built as an ensemble of four ML modules: a \textbf{Diagnostic} module (clustering $\rightarrow$ knowledge priors), a \textbf{Behavioural Diagnosis} neural network (response metadata $\rightarrow$ a behavioural profile $\rightarrow$ a treatment plan), a \textbf{Question Bandit} (LinUCB $\rightarrow$ next question), and a \textbf{Progress-Tracking} module (Bayesian Knowledge Tracing $\rightarrow$ mastery). This note states the project's \textbf{main research question and its sub-questions}, enumerates \textbf{all student cases from the learner's perspective} with their observable signals and treatments, lists the \textbf{log-derived metrics/insights} the system must compute, and defines a concrete, cold-start-safe \textbf{question-difficulty formula}. The Chinese-math \texttt{XES3G5M} dataset used during prototyping is a \emph{proxy only}: no Egyptian-science behavioural dataset exists yet, which makes cold-start and domain transfer first-class research problems.
\end{abstract}

% =====================================================================
\section{Context and Scope}
\textbf{Domain:} Egyptian national-curriculum \textbf{science} (physics, chemistry, biology), Arabic language, roughly grades 4--12.\\
\textbf{Product:} a competitive, timed quiz \emph{game}; every question carries rich metadata (response time, chosen distractor, question type, cognitive load).\\
\textbf{Goal:} \emph{personalised sequencing} --- recommend the next question/level per student to maximise learning gain, subject to keeping the learner inside their Zone of Proximal Development and respecting behavioural ``treatment'' constraints.\\
\textbf{Data status:} the pipeline was prototyped on \texttt{XES3G5M} (Chinese elementary math: $\sim$5.1M interactions, 14{,}453 students, 7{,}618 questions, 865 knowledge components, mean error rate $\approx0.20$). This is a \textbf{methodology proxy}, not Noxed's data. The real target data (Egyptian science, with behaviour labels) \textbf{does not exist yet} and must be instrumented and/or synthesised --- see \S\ref{sec:data}.

% =====================================================================
\section{Main Research Question and Sub-Questions}\label{sec:rq}

\rqbox{Main RQ.}{Can we infer, from a student's in-game response data alone --- correctness patterns, solving time, and mistake structure --- both (i) \emph{what they do not know} (knowledge gaps) and (ii) \emph{how they behave under test conditions} (stable behavioural traits), and use these together to \textbf{recommend the next question / difficulty level} that maximises the student's learning gain?}

\medskip
This decomposes into one sub-question per module plus two cross-cutting ones.

\subsection{RQ1 --- Knowledge diagnosis (Diagnostic module)}
\begin{itemize}
\item[\textbf{RQ1}] From the knowledge components (KCs) present in a student's \emph{wrong} vs.\ \emph{correct} answers, can hierarchical clustering recover their weak vs.\ strong topics and yield useful initial mastery priors?
\item[RQ1.1] Which features (KC-route embedding, tree depth, co-occurrence, error rate) make wrong-answer and correct-answer clusters \emph{meaningful} rather than arbitrary?
\item[RQ1.2] How should cluster \emph{size} and cross-cluster \emph{multiplicity} map to a prior in $[0,1]$ (larger/repeated wrong-answer clusters $\Rightarrow$ lower prior; the inverse for correct)?
\item[RQ1.3] How do we exploit the KC \emph{prerequisite tree} so a low prior on a parent KC suppresses its children?
\end{itemize}

\subsection{RQ2 --- Behavioural diagnosis (Behavioural NN)}
\begin{itemize}
\item[\textbf{RQ2}] Can a neural network classify a student's \emph{behavioural scenario} (e.g.\ poor time management, impulsivity, test anxiety, weak context-switching) from response metadata, and how well does each latent trait separate from the others?
\item[RQ2.1] \emph{Which observable metrics operationalise each trait?} (See \S\ref{sec:metrics}.)
\item[RQ2.2] Are these traits \emph{stable} enough (across matches/sessions) to be treated as latent student properties rather than momentary noise?
\item[RQ2.3] What is the \emph{ground truth}? Absent labels, can synthetic ``scenario'' generators and/or weak-supervision heuristics (wheel-spinning, gaming, carelessness rules) train the NN, and does it transfer to real students?
\item[RQ2.4] Does near-identical behaviour reliably map to near-identical treatment plans (smoothness/consistency of the score$\rightarrow$plan rule layer)?
\end{itemize}

\subsection{RQ3 --- Next-question recommendation (Question Bandit)}
\begin{itemize}
\item[\textbf{RQ3}] Given a student's priors and treatment plan, can a contextual bandit (LinUCB) select the next question that maximises expected learning gain while staying inside the ZPD?
\item[RQ3.1] What is the right \emph{reward}? The current form $R=\tfrac{1}{\sum \text{priors}}\cdot Q_{\text{diff}}\cdot Q_{\text{learn}}$ favours low-mastery, harder, high-improvement items --- does it need a \emph{desirable-difficulty} term so success probability stays near a target ($\approx$\,0.7--0.85)? (See \S\ref{sec:reward}.)
\item[RQ3.2] What context features (student ability vector, item features, behavioural scores) should be the LinUCB arm context?
\item[RQ3.3] How do we \emph{validate} a bandit policy offline (off-policy evaluation) before any live student is exposed to it?
\end{itemize}

\subsection{RQ4 --- Progress \& mastery (Progress Tracking)}
\begin{itemize}
\item[\textbf{RQ4}] Does BKT-based mastery tracking, updated after each match, correctly detect when a student has mastered a level and should be advanced?
\item[RQ4.1] How are BKT parameters $P(L_0),P(T),P(S),P(G)$ initialised from static question features before response data exists? (See \S\ref{sec:priors}.)
\item[RQ4.2] Does solving a question actually reduce a student's error on its KCs (practice effect), and by how much (learning $\Delta$)?
\end{itemize}

\subsection{RQ5 \& RQ6 --- Cross-cutting}
\begin{itemize}
\item[\textbf{RQ5 (Cold start / transfer)}] With no Egyptian-science behavioural data, how far can models trained on proxy datasets (\texttt{XES3G5M}, ASSISTments, EdNet) transfer, and what is the minimum data Noxed must collect itself?
\item[\textbf{RQ6 (Efficacy / ATI)}] Does trait-matched sequencing (\emph{aptitude--treatment interaction}) actually improve learning versus a difficulty-only baseline, and for which student profiles?
\end{itemize}

% =====================================================================
\section{All Student Cases (Learner's Perspective)}\label{sec:cases}
The behavioural target space is a four-family taxonomy of student ``scenarios'' (authored for the Egyptian context). Each case is defined by an \textbf{observable in-game signal} (what the NN/rules detect) and a \textbf{treatment} (a constraint added to the treatment plan). These are the classes the Behavioural NN must ultimately recognise.

\subsection{Family A --- Content Gaps (academic mastery)}
\begin{tabularx}{\textwidth}{@{}p{3.6cm}YY@{}}
\toprule
\textbf{Case} & \textbf{Diagnosis (observable signal)} & \textbf{Treatment (plan constraint)}\\
\midrule
Knowledge Absence & Random guessing on ordering/matching items; very high failure on a lesson; distractors chosen with no scientific logic. & Stop complex application items; serve True/False and direct definitions to rebuild the base.\\
Weak Prior Knowledge & Wrong answers rooted in \emph{foundational} errors (e.g.\ an old formula), with long time spent. & Insert scaffolding/prerequisite items before the target KC; give rule-reminder hints.\\
Rote w/o Understanding & Matches a term to its literal definition correctly, but fails to \emph{apply} it; swaps two adjacent steps in ordering items. & Switch the same content to visual item types (Image MCQ / drag-and-drop) to bind concept to image, not text.\\
Misconceptions & Places elements in logic-reversed order; picks a distractor built on a common science myth. & ``Cognitive-conflict'' cue (distinct SFX) + refutational feedback that dismantles the wrong model.\\
\bottomrule
\end{tabularx}

\subsection{Family B --- Cognitive Skills (information processing)}
\begin{tabularx}{\textwidth}{@{}p{3.6cm}YY@{}}
\toprule
\textbf{Case} & \textbf{Diagnosis (observable signal)} & \textbf{Treatment (plan constraint)}\\
\midrule
Working Memory & Performance collapses on high Cognitive-Load-Index items (many variables / steps). & Fade out background music during the item to free focus; prompt ``use pen and paper''.\\
Processing Speed & Correct but always consumes maximum time; more errors on time-pressure items. & Reduce time-pressure items early; nudge ``accuracy first, then speed''.\\
Attention (selective/sustained) & Performance drops on the \emph{last} items of a match; errors when a visual distractor (fog/monster) is present. & 1-second silence-pause to reset attention before a distractor item; avoid linguistically complex items late in the game.\\
Cognitive Flexibility & Sharp drop / slower responses right after a \emph{topic switch} (context switch cost). & Progress from blocked practice to gradual interleaving, with an advance warning of a rule change.\\
\bottomrule
\end{tabularx}

\subsection{Family C --- Exam Skills (tactical)}
\begin{tabularx}{\textwidth}{@{}p{3.6cm}YY@{}}
\toprule
\textbf{Case} & \textbf{Diagnosis (observable signal)} & \textbf{Treatment (plan constraint)}\\
\midrule
Question Interpretation & Picks the ``obvious'' wrong answer on negation/exception (\emph{except}, \emph{not}) items, and does so \emph{fast}. & Fading: attach a distinctive cue to negation items to build attention, then gradually withdraw it.\\
Time Mgmt.\ \& Perfectionism & Spends 90\% of time on trivially easy items; or times out on many items. & Internal-pacer SFX at 50\% time elapsed; train with short, fast items to break hesitation.\\
\bottomrule
\end{tabularx}

\subsection{Family D --- Psychological / Behavioural (affect under pressure)}
\begin{tabularx}{\textwidth}{@{}p{3.6cm}YY@{}}
\toprule
\textbf{Case} & \textbf{Diagnosis (observable signal)} & \textbf{Treatment (plan constraint)}\\
\midrule
Impulsiveness & Responds in \textless\,30\% of the theoretical time and picks the obvious/impulsive distractor. & Slow the music tempo; disable auto score-multipliers to force deliberation.\\
Stress \& Overthinking & Gross errors on easy items, or doubled time, \emph{only} under stress triggers (last place, direct threat from a rival). & Graded exposure: heartbeat-tempo music in critical moments to train focus under pressure; give a confidence-restoring direct item next round.\\
\bottomrule
\end{tabularx}

\medskip
\noindent\textit{Note: these treatments are game-side interventions (audio, item-type, pacing, scaffolding); the Behavioural module emits them only from a fixed \textbf{treatment vocabulary} so the bandit's filtering is well-defined.}

% =====================================================================
\section{Metrics \& Insights the System Must Compute}\label{sec:metrics}
To detect the cases above, every response is reduced to log-derived features. These are grouped by the construct they indicate; they are the candidate inputs to the Behavioural NN and the signals behind each treatment trigger.

\begin{tabularx}{\textwidth}{@{}p{4.9cm}p{3.5cm}Y@{}}
\toprule
\textbf{Metric (log-derived)} & \textbf{Construct} & \textbf{Definition / signal}\\
\midrule
Normalised response time $\tau=t/t_{\text{expected}}$ & Speed, impulsivity & Ratio of solving time to the item's expected time.\\
Fast-error rate & Impulsivity & Fraction of wrong answers with $\tau<0.3$.\\
Slow-correct rate & Processing speed & Fraction of correct answers at $\tau\ge1$ (max time).\\
Timeout rate & Time management & Fraction of items with no answer before the timer.\\
Time-on-easy share & Perfectionism & Time spent on lowest-difficulty items / total time.\\
Carelessness (slip) & Attention & Errors on items whose KC mastery $P(L)$ is high.\\
Wheel-spinning & Persistence & $\ge N$ attempts on a KC with no rise in mastery.\\
Gaming-the-system & Disengagement & Systematic fast wrong answers / hint abuse.\\
Distractor-type profile & Misconception & Distribution over impulsive vs.\ misconception vs.\ prerequisite distractors chosen.\\
Negation-item error gap & Interpretation & Error on negation/``except'' items minus error on plain items.\\
Context-switch cost & Cognitive flexibility & $\Delta$(error, RT) on the first item after a topic change vs.\ within-topic.\\
Cognitive-load sensitivity & Working memory & Slope of error vs.\ Cognitive-Load-Index.\\
End-of-match decay & Sustained attention & Error on last-$k$ items minus first-$k$ items.\\
Pressure sensitivity & Test anxiety & Error/RT change under stress game-states (last place, rival threat).\\
Recovery rate & Emotional regulation & Accuracy on the item following a wrong/stress event.\\
Learning $\Delta$ per KC & Practice effect & Error-rate drop on a KC before vs.\ after solving related items.\\
Differentiation $\delta$ & Item quality & Item error minus the student's baseline error on its KCs.\\
Attempt count / KC & Persistence & Number of tries before first success on a KC.\\
Answer-change rate & Anxiety / overthinking & Frequency of switching a selected answer before submit (if UI allows).\\
Session engagement & Affect & Match completion, drop-off, re-entry latency.\\
\bottomrule
\end{tabularx}

\medskip
\noindent\textbf{Key project insights.}
\begin{enumerate}[leftmargin=1.4em]
\item \textbf{Desirable difficulty, not maximum difficulty.} Learning is fastest when success probability sits near $0.7$--$0.85$; the reward must target that band, not simply the hardest item (\S\ref{sec:reward}).
\item \textbf{Behaviour needs its own labels.} Correctness alone cannot label anxiety or impulsivity; ground truth must come from synthetic scenarios and/or weak-supervision rules --- this is the single biggest data risk.
\item \textbf{Cold start dominates the roadmap.} With no Egyptian-science data, item difficulty and BKT priors must be derived from \emph{static} features first, then blended with evidence as it accrues (\S\ref{sec:difficulty}).
\item \textbf{The KC tree is the backbone.} Prerequisite structure (up to depth 7 in the proxy) constrains both what to recommend next and how errors propagate to parent topics.
\item \textbf{Practice effect is measurable and directional.} A negative before/after error delta on a KC is direct evidence the sequence is teaching --- use it to validate RQ4.
\end{enumerate}

\noindent\textbf{Measurement caveats.} (1) Almost every raw metric must be \emph{item-difficulty-normalised} (per-item $z$-score / lognormal residual) before it separates \emph{behaviour} from \emph{ability}. (2) Under a competitive timer, test anxiety, impulsivity and rapid-guessing are mutually confounded --- prefer \emph{within-student} contrasts (same student across pressure/difficulty conditions) over raw cross-sectional values. (3) The ${\approx}0.8$ target success probability is a design heuristic, not a fixed constant, and ATI trait$\rightarrow$treatment matches must be A/B-validated, not assumed (Cronbach \& Snow's own caveat).

% =====================================================================
\section{Question-Difficulty Formula}\label{sec:difficulty}
Difficulty $D\in[0,1]$ is required by the bandit reward and by the BKT priors, but at launch \emph{no response data exists} for Egyptian-science items. We therefore define difficulty in three layers: a \textbf{static} estimate from authoring features (works at cold start), an \textbf{empirical} estimate from response data (works once logs exist), and a \textbf{confidence-weighted blend} that transitions smoothly from one to the other.

\subsection{Static (cold-start) difficulty}
Let each item carry expert-authored features, each normalised to $[0,1]$:
\[
\begin{array}{llll}
x_1 & \text{cognitive-load index (steps, variables)} & x_4 & \text{item-type base (TF}<\text{MCQ}<\text{multi}<\text{fill/DnD)}\\
x_2 & \text{KC depth in the tree }(\text{depth}/7) & x_5 & \text{linguistic challenge (negation/``except'')}\\
x_3 & \text{KC integration }(\min(k,K)/K) & x_6 & \text{distractor plausibility}
\end{array}
\]
Combine them through a logistic link so the output is \emph{bounded in $[0,1]$ by construction}, not by clipping:
\[
\boxed{\,D_{\text{stat}} \;=\; \sigma\!\Big(\beta_0+\textstyle\sum_{i=1}^{6}\beta_i x_i\Big),\qquad \sigma(z)=\frac{1}{1+e^{-z}}\,}
\]
with non-negative weights $\beta_i$ (higher load / depth / integration / type / language / distractor strength $\Rightarrow$ harder). Weights are set by expert calibration first, then fit by logistic regression against $D_{\text{emp}}$ once data arrives.

\subsection{Empirical difficulty (once responses exist)}
Raw error rate is unreliable for rarely-attempted items, so shrink it toward the global mean $\mu$ (Beta--Binomial / empirical Bayes) with pseudo-count $m$:
\[
\boxed{\,D_{\text{emp}} \;=\; \frac{w+\mu\,m}{a+m}\,}\qquad
\begin{array}{l}
a=\text{\#attempts},\; w=\text{\#wrong},\\
\mu=\text{global mean error }(\approx0.20\text{ in the proxy}),\; m\approx20 .
\end{array}
\]
Optionally correct for \emph{intrinsic} hardness beyond topic mastery using the differentiation delta $\delta$ (item error minus the student's baseline error on the item's KCs):
\[
D_{\text{emp}}^{\star}=\operatorname{clip}\!\big(D_{\text{emp}}+\gamma\,\delta,\,0,\,1\big),\qquad \gamma\in[0,1].
\]

\subsection{Blended difficulty (cold start $\rightarrow$ data-driven)}
Weight the empirical term by how much evidence backs it, $\lambda(a)=a/(a+a_0)$ (reliability constant $a_0\approx30$):
\[
\boxed{\,D \;=\; \lambda(a)\,D_{\text{emp}}^{\star} \;+\; \big(1-\lambda(a)\big)\,D_{\text{stat}}\,}
\]
At $a=0$ (a brand-new Egyptian-science item) $D=D_{\text{stat}}$; as responses accumulate $D\to D_{\text{emp}}^{\star}$. Being a convex combination of two quantities in $[0,1]$, $D\in[0,1]$ automatically --- a \textbf{property directly testable with Hypothesis} for any input.

\subsection{From difficulty to BKT priors}\label{sec:priors}
The same $D$ (and load $x_1$, option count $o$, type) seeds the per-item BKT parameters, each clipped to identifiable ranges ($P(S),P(G)<0.5$):
\[
P(G)=\operatorname{clip}\!\big(g_0\,(1-D)\cdot \tfrac{1}{o},\,0,\,0.3\big),\quad
P(S)=\operatorname{clip}\!\big(s_0\,(0.5+0.5D)(0.5+0.5x_1),\,0,\,0.3\big),
\]
\[
P(L_0)=\operatorname{clip}\!\big(0.5(1-D),\,0,\,1\big)\ \text{(fallback; the Diagnostic module overrides it)},\quad
P(T)=t_0\,(1-0.5D).
\]

\subsection{Difficulty in the bandit reward}\label{sec:reward}
The given reward $R=\dfrac{1}{\sum_j \text{prior}_j}\;Q_{\text{diff}}\;Q_{\text{learn}}$ rewards low mastery, high difficulty and high learn-rate. To keep the learner in the ZPD, replace raw $Q_{\text{diff}}$ with a \emph{desirable-difficulty match} on predicted success probability
\(p=P(L)\,(1-P(S))+(1-P(L))\,P(G)\):
\[
\text{Match}=\exp\!\Big(-\tfrac{(p-p^\star)^2}{2\tau^2}\Big),\quad p^\star\approx0.8,\qquad
R' = \frac{1}{\sum_j \text{prior}_j}\;\text{Match}\;Q_{\text{learn}} .
\]
This peaks when the item is challenging-but-winnable rather than merely hard.

% =====================================================================
\section{Validation Plan (properties to test)}
\begin{itemize}
\item \textbf{Partitioning:} every clustered answer belongs to exactly one cluster.
\item \textbf{Ranges:} behavioural scores, $P(L)$, and every $D$ variant $\in[0,1]$; $\Delta$mastery $\in[-1,1]$ (fuzz with Hypothesis).
\item \textbf{Vocabulary:} every emitted treatment constraint is a member of the fixed treatment vocabulary.
\item \textbf{Smoothness:} near-identical behavioural score vectors map to the same/very similar treatment plan.
\item \textbf{Bandit safety:} the selected question is always inside the constraint-filtered candidate set.
\item \textbf{Efficacy:} offline off-policy evaluation shows the policy beats a difficulty-only baseline before any live rollout (RQ6).
\end{itemize}

% =====================================================================
\section{Data Strategy: Proxy, Cold Start, Synthesis}\label{sec:data}
A verified survey of public datasets yields \textbf{three non-overlapping buckets}: (a) large knowledge-tracing logs with rich per-attempt metadata but \emph{no} behavioural labels; (b) small human-observed affect/behaviour corpora (BROMP lineage) that are log-linked but tiny; (c) webcam/video affect corpora with labels but \emph{no} response metadata. \textbf{No public dataset offers both scale and behavioural ground truth}, and \textbf{none matches Noxed's target} (Arabic $\times$ school-science $\times$ per-attempt logs).

\smallskip
\noindent\textbf{Proxy / transfer datasets (methodology only).} RT = response time, Att = attempts, Hint = hint use, KC = skill tags, Aff = affect labels.

\par\nopagebreak
\begin{tabularx}{\textwidth}{@{}p{3.2cm}Y ccccc@{}}
\toprule
\textbf{Dataset} & \textbf{Domain / size} & RT & Att & Hint & KC & Aff \\
\midrule
\texttt{XES3G5M} (current) & CN elem.\ math; 18{,}066 stu, 7{,}652 q, 865 KC, 5.55M int & $\sim$ & -- & -- & \checkmark\checkmark & -- \\
ASSISTments 2012-13 \emph{+affect} & US math; 27K stu, 2.5M int & \checkmark & \checkmark & \checkmark & \checkmark & (pred.) \\
ASSISTments 2017 & US math; clickstream, 942K rows & \checkmark & \checkmark & \checkmark & \checkmark & (pred.) \\
EdNet KT1--KT4 (Riiid) & EN TOEIC; 784K stu, 131M int & \checkmark\checkmark & \checkmark & $\sim$ & \checkmark & -- \\
Eedi / NeurIPS-2020 & UK math MCQ; $>$20M answers & -- & \checkmark & -- & \checkmark\checkmark & -- \\
Junyi 2020 (2015) & TW math; 72K stu, 16M att & \checkmark & \checkmark & \checkmark & \checkmark & -- \\
KDD Cup 2010 & US algebra; $\sim$30M steps & \checkmark & \checkmark & \checkmark & \checkmark & -- \\
Squirrel AI ElemMath2021 & CN elem.\ math; 125K stu, difficulty+prereqs & \checkmark & \checkmark & -- & \checkmark & -- \\
\bottomrule
\end{tabularx}

\smallskip
\noindent\textbf{Closest to Noxed's domain: TIMSS Grade 4/8 \emph{science}} (IEA) --- \textbf{Egypt participates}, many Arabic countries included, item-level responses public, and eTIMSS adds response-time process data. It is cross-sectional assessment data (no hints/retries/sequences), so mine it for \emph{science-item content and difficulty realism}, and use a math KT log (EdNet / Junyi / ASSISTments) as the \emph{format \& behaviour} proxy. (PISA science is computer-based with time-on-task but Egypt does not take PISA.)

\smallskip
\noindent\textbf{Cold-start reality.} There is no Arabic, per-attempt, school-science tutoring dataset publicly available; moreover \emph{test anxiety} and \emph{impulsiveness} are not even standard observation categories (they come from self-report). Noxed must therefore:
\begin{enumerate}[leftmargin=1.4em,itemsep=1pt]
\item \textbf{Instrument its own game} from day one --- log response time, attempts, chosen distractor, lifeline/skip/answer-change events, and stress game-states.
\item \textbf{Synthesise labelled behavioural scenarios} --- parametrised ``impulsive'' (short RT $+$ lure-distractor), ``anxious'' (RT inflation $+$ decay under pressure), ``slow-processor'', ``gamer'' (implausibly short RT $+$ wrong) students --- to pre-train the Behavioural NN.
\item \textbf{Anchor weak supervision} on ASSISTments-with-affect ``silver'' labels and the BROMP/sensor-free-detector definitions, then validate a subset against a short Arabic-validated self-report (e.g.\ Cognitive Test Anxiety Scale).
\end{enumerate}
Item difficulty and BKT priors come from the \emph{static} formula (\S\ref{sec:difficulty}) until enough real responses accrue.

% =====================================================================
\section{Reusable Open-Source Landscape}
\textbf{Knowledge tracing / mastery:} \texttt{pyBKT} (BKT priors \& mastery --- already in stack), \texttt{pyKT} and \texttt{EduKTM} (DKT/DKVMN/SAKT/AKT baselines), \texttt{py-fsrs}/\texttt{ebisu} (forgetting/review timing).\\
\textbf{Adaptive tutor reference:} \texttt{OATutor} (open BKT-based ITS with a mastery$\rightarrow$selection loop --- closest architectural analog).\\
\textbf{Contextual bandit (M3):} \texttt{contextualbandits} (David Cortes --- LinUCB \emph{plus} off-policy evaluation), \texttt{Vowpal Wabbit} (scale), \texttt{catsim}/\texttt{EduCAT} (information-theoretic CAT as a principled benchmark).\\
\textbf{Item calibration:} \texttt{py-irt}, \texttt{girth}, \texttt{mirt} (difficulty/discrimination to seed the difficulty formula \& priors).\\
\textbf{Behavioural (M2):} no turnkey library --- reuse the \emph{definitions} of wheel-spinning / gaming / carelessness detectors (Baker, Paquette, Beck) and \texttt{pyKT} auxiliary-task heads; train a custom PyTorch classifier. \emph{(Note: ``OATutor'', from the same lab as pyBKT, is the closest end-to-end open-source analog to Noxed's mastery$\rightarrow$selection loop.)}

% =====================================================================
\section*{References}
\small
\begin{enumerate}[leftmargin=1.5em,itemsep=0pt]
\item Corbett, A.T.\ \& Anderson, J.R.\ (1995). Knowledge Tracing: Modeling the acquisition of procedural knowledge. \emph{UMUAI} 4(4), 253--278.
\item Piech, C.\ et al.\ (2015). Deep Knowledge Tracing. \emph{NeurIPS} 28.
\item Zhang, J.\ et al.\ (2017). Dynamic Key-Value Memory Networks for Knowledge Tracing. \emph{WWW}, 765--774.
\item Ghosh, A., Heffernan, N.\ \& Lan, A.S.\ (2020). Context-Aware Attentive Knowledge Tracing (AKT). \emph{KDD}, 2330--2339.
\item Baker, R.S.\ et al.\ (2008). Developing a generalizable detector of when students game the system. \emph{UMUAI} 18(3), 287--314.
\item Beck, J.E.\ \& Gong, Y.\ (2013). Wheel-Spinning: students who fail to master a skill. \emph{AIED}, LNCS 7926, 431--440.
\item San Pedro, M.O.Z., Baker, R.S.\ \& Rodrigo, M.M.T.\ (2011). Detecting carelessness through contextual estimation of slip probabilities. \emph{AIED}, 304--311.
\item D'Mello, S.\ \& Graesser, A.\ (2012). Dynamics of affective states during complex learning. \emph{Learning \& Instruction} 22(2), 145--157.
\item Ocumpaugh, J., Baker, R.S.\ \& Rodrigo, M.M.T.\ (2015). \emph{BROMP 2.0 Training Manual}.
\item Cronbach, L.J.\ \& Snow, R.E.\ (1977). \emph{Aptitudes and Instructional Methods} (ATI). Irvington.
\item Miyake, A.\ et al.\ (2000). The unity and diversity of executive functions. \emph{Cognitive Psychology} 41(1), 49--100. (See also Diamond, 2013, \emph{Annu.\ Rev.\ Psychol.} 64.)
\item Cassady, J.C.\ \& Johnson, R.E.\ (2002). Cognitive test anxiety and academic performance. \emph{Contemp.\ Educ.\ Psychol.} 27(2), 270--295.
\item Wise, S.L.\ \& Kong, X.\ (2005). Response-time effort: a measure of examinee motivation. \emph{Applied Measurement in Education} 18(2), 163--183.
\item van der Linden, W.J.\ (2006/2007). Lognormal / hierarchical models of response times. \emph{JEBS} 31(2); \emph{Psychometrika} 72(3).
\item Clement, B.\ et al.\ (2015). Multi-armed bandits for intelligent tutoring systems (ZPDES). \emph{JEDM} 7(2), 20--48.
\item Settles, B.\ \& Meeder, B.\ (2016). A trainable spaced-repetition model (Half-Life Regression). \emph{ACL}, 1848--1858.
\item Doroudi, S., Aleven, V.\ \& Brunskill, E.\ (2019). Where's the reward? RL for instructional sequencing. \emph{IJAIED} 29(4), 568--620.
\item Bloom, B.S.\ (1984). The 2-sigma problem. \emph{Educational Researcher} 13(6), 4--16. \; Bjork, R.A.\ (1994). Desirable difficulties.
\item Liu, Z.\ et al.\ (2023). \texttt{XES3G5M}: a knowledge-tracing benchmark with auxiliary information. \emph{NeurIPS Datasets \& Benchmarks}.
\end{enumerate}

\end{document}
